\section{Fourier 分析}

\subsection*{Fourier Series}

Origin: Heat conduction

$T$: $\mathbb{R}^3\to\mathbb{R}$ temperature

For any domain $\Omega$

\begin{align*}
    \frac{\mathrm{d}}{\mathrm{d}t}\int_\Omega T\left(\overrightarrow{r},t\right)\,\mathrm{d}V\left(\overrightarrow{r}\right)
     & =\oiint_{\partial\Omega}\nabla T\left(\overrightarrow{r},t\right)\cdot\mathrm{d}\overrightarrow{S}      \\
     & =\int_{\Omega}\partial_t T\left(\overrightarrow{r},t\right)\,\mathrm{d}V\left(\overrightarrow{r}\right) \\
     & =\int_\Omega\Delta T\left(\overrightarrow{r},t\right)\,\mathrm{d}V\left(\overrightarrow{r}\right)
\end{align*}

$\partial_t T=k^2\Delta T$

\begin{itemize}
    \item Assuming $k\equiv 1$, $\partial_t T=\Delta T$
    \item $n=1$, $\begin{cases}
                  \partial_t T=\partial_{xx}T \\
                  T(x,0)=f(x)     & x\in[0,l] \\
                  T(0,t)=T(l,t)=0 & \forall t
              \end{cases}$
    \item \allbold{boundary value problem on $[0,l]\times[0,+\infty)$}
\end{itemize}

\paragraph*{Heat equation 热(传导方程)}

$k$: conducting rate

Step 1: Looking for solution of separating variables, i.e. $T(x,t)=\varphi(x)\psi(t)$

$\implies\varphi(x)\psi'(t)=\varphi''(x)\psi(t)$

$\implies\frac{\varphi''(x)}{\varphi(x)}=\frac{\psi'(t)}{\psi(t)}\equiv\lambda$ csonstant

Solution: $\psi(t)=\psi(0)\mathrm{e}^{\lambda t}$

$\varphi''(x)-\lambda\varphi(x)=0\leftarrow$ 2nd-order ODE of constant coefficients

$\begin{cases}
        \varphi(x)=c_1\mathrm{e}^{\sqrt{\lambda}x}+c_2\mathrm{e}^{-\sqrt{\lambda}x} & \lambda>0 \\
        \varphi(x)=c_1\cos\sqrt{-\lambda}x+c_2\sin\sqrt{-\lambda}x                  & \lambda<0 \\
    \end{cases}$

$x<0$, $\left(\frac{\mathrm{d}}{\mathrm{d}x}+i\sqrt{-\lambda}\right)\underset{\Phi(x)}{\left(\frac{\mathrm{d}}{\mathrm{d}x}-i\sqrt{-\lambda}\right)\varphi(x)}=0$

$\left(\mathrm{e}^{i\sqrt{-\lambda}x}\Phi(x)\right)'=0$

$\implies\Phi(x)=-\varphi'(x)-i\sqrt{-\lambda}\varphi(x)=c_i\mathrm{e}^{-i\sqrt{-\lambda}}x$

$\left(\mathrm{e}^{-i\sqrt{-\lambda}x}\varphi(x)\right)'=\tilde{c_1}\mathrm{e}^{-2i\sqrt{-\lambda}}x+\tilde{c_2}\mathrm{e}^{-i\sqrt{-\lambda}x}$, $\tilde{c_1},\tilde{c_2}\in\mathbb{C}$

$\left(\frac{\mathrm{d}}{\mathrm{d}x}-\lambda_1\right)\cdot\left(\frac{\mathrm{d}}{\mathrm{d}x}-\lambda_n\right)\varphi(x)=\varphi^{(n)}+a_1\varphi^{(n-1)}+\dots+a_n\varphi(x)=f(x)$, $a_i$ constant

$\impliedby\lambda^n+a_1\lambda^{n-1}+\dots+a_n=(\lambda-\lambda_1)\cdots(\lambda-\lambda_n)$

$\varphi''(x)-\lambda\varphi(x)=\left(\frac{\mathrm{d}}{\mathrm{d}x}+\sqrt{\lambda}\right)\left(\frac{\mathrm{d}}{\mathrm{d}x}-\sqrt{\lambda}\right)\varphi(x)=0$

$\Phi'(x)+\sqrt{\lambda}\Phi(x)=0$

$\left(\mathrm{e}^{\sqrt{\lambda}x}\Phi(x)\right)'=\mathrm{e}^{\sqrt{\lambda}x}\left(\Phi'(x)+\sqrt{\lambda}\Phi(x)\right)=0$

$\implies\Phi(x)=\varphi'(x)-\sqrt{\lambda}\varphi(x)=c_1\mathrm{e}^{-\sqrt{\lambda}x}$

$\mathrm{e}^{-\sqrt{\lambda}x}\left(\varphi'(x)-\sqrt{\lambda}\varphi(x)\right)=\left(\mathrm{e}^{-\sqrt{\lambda}x}\varphi(x)\right)'=c_1\mathrm{e}^{-2\sqrt{\lambda}x}$

$\implies\mathrm{e}^{-\sqrt{\lambda}x}\varphi(x)=\tilde{c_1}\mathrm{e}^{-2\sqrt{\lambda}x}+c_2$

Notice that the zero boundary condition rules out $\lambda>0$

Take $\lambda<0$, for convenience, replace $\lambda$ by $-\lambda^2$

$\varphi''(x)+\lambda^2\varphi(x)=0$

$\varphi(x)=c_1\cos\lambda x+c_2\sin\lambda x$

$\begin{cases}
        \varphi(0)=c_1=0 \\
        \varphi(l)=c_1\cos\lambda l+c_2\sin\lambda l=0
    \end{cases}\implies\begin{cases}
        \lambda l=n\pi\text{, for }n=1,2,\dots \\
        \lambda-\lambda_n=\frac{n\pi}{l}
    \end{cases}$

$c_1\mathrm{e}^{\sqrt{\lambda}x}+c_2\mathrm{e}^{-\sqrt{\lambda}x}=0$, $x=0,l$

$\begin{cases}
        c_1+c_2=0 \\
        c_1\mathrm{e}^{\sqrt{\lambda}l}+c_2\mathrm{e}^{-\sqrt{\lambda}l}=0
    \end{cases}\implies c_1=c_2=0$

$$\boxed{T_n(x,t)}=\mathrm{e}^{-\left(\frac{n\pi}{l}\right)^2}\sin\frac{n\pi}{l}x$$

Since the heat is linear,
$$\sum_{n=1}^{\infty}a_nT_n(x,t)=\sum_{n=1}^{\infty}a_n\mathrm{e}^{-\left(\frac{n\pi}{l}\right)^2t}\sin\frac{n\pi}{l}x$$
is also linear

Plug in $t=0$, $\sum_{n=1}^{\infty}a_n\sin\frac{n\pi}{l}x=f(x)$ (for any $f:\,[0,l]\to\mathbb{R}$ with $f(0)=f(l)=0$)

$\sin\frac{\pi}{l}x,\sin\frac{2\pi}{l}x,\sin\frac{3\pi}{l}x,\dots$

In general, $f(x)\xlongequal{?}\frac{a_0}{2}+\sum_{n=1}^{\infty}\left(a_n\cos nx+b_n\sin nx\right)$

\begin{definition}
    Let $f:\,\mathbb{R}\to\mathbb{R}$, if $\exists T>0$ s.t. $f(x+T)=f(x)$ for any $x\in\mathbb{R}$, then $T$ is called ``a period'' of $f$, ($2T,3T,4T,\dots$ are all periods of $f$)

    (If such $T$ exists, $f$ is called periodoc)

    If there exists a minimum period, it is called the minimal period
\end{definition}

$P_f\triangleq\left\{T>0\middle|T\text{is a period of }f\right\}\implies\inf P_f\geq 0$

\begin{example}
    Given $n=1,2,\dots$, $\sin nx\leadsto$ minimal period $\frac{2\pi}{n}$
\end{example}

\begin{example}
    $D(x)\leadsto$ It has no minimal period
\end{example}

If $l$ is a period of $f$, then $g(x)\triangleq f\left(\frac{l}{2\pi}x\right)$ has $2\pi$ as its period

\begin{definition}[Fourier Series]
    Let $f:\,\mathbb{R}\to\mathbb{R}$ be a Riemann integrable function and suppose it has period $2\pi$, then
    $$\frac{a_0}{2}+\sum_{n=1}^{\infty}\left(a_n\cos nx+b_n\sin nx\right)$$
    is called the Fourier series of $f$, where
    $$\begin{cases}
            a_n=\frac{1}{\pi}\int_{-\pi}^{\pi}f(x)\cos nx\,\mathrm{d}x=\left<f,\cos nx\right>_{L^2} \\
            b_n=\frac{1}{\pi}\int_{-\pi}^{\pi}f(x)\sin nx\,\mathrm{d}x=\left<f,\sin nx\right>_{L^2}
        \end{cases}\quad n=0,1,2,\dots$$
    are called the Fourier coefficients of $f$
\end{definition}

\begin{remark}
    Changing the values of $f$ at finitely many points does not change its Fourier series
\end{remark}

\subparagraph*{Fact}

\begin{align*}
    \frac{1}{\pi}\int_{-\pi}^{\pi}\cos nx\cdot\cos mx\,\underset{=\left<\cos nx,\cos mx\right>_{L^2}}{\mathrm{d}x} & =\delta_{mn} \\
    \frac{1}{\pi}\int_{-\pi}^{\pi}\sin nx\cdot\sin mx\,\underset{=\left<\sin nx,\sin mx\right>_{L^2}}{\mathrm{d}x} & =\delta_{mn}
\end{align*}

$\implies\left\{\frac{1}{\sqrt{2}},\cos x,\sin x,\cos 2x,\sin 2x,\dots\right\}$ gives an orthonormal set of vectors in $R[-\pi,\pi]$

$\implies$ If we introduce the notation:

$$\left<f,g\right>_{L^2}\triangleq\frac{1}{\pi}\int_{-\pi}^{\pi}f(x)g(x)\,\mathrm{d}x,\quad\forall f,g\in R[-\pi,\pi]$$

\begin{itemize}
    \item Bilinear
    \item Symmetric
    \item Positive definite? ($f\sim g$ if $f=g$ almost everywhere)
\end{itemize}

\allbold{$\implies$ Complete metric space?}

The Fourier series of $f$ is actually
$$\left<f,\frac{1}{\sqrt{2}}\right>_{L^2}\frac{1}{\sqrt{2}}+\left<f,\varphi_1\right>_{L^2}\varphi_1(x)+\left<f,\psi_1\right>_{L^2}\psi_1(x)+\left<f,\varphi_2\right>_{L^2}\varphi_2(x)+\left<f,\psi_2\right>_{L^2}\psi_2(x)+\dots$$
where $\begin{cases}
        \varphi_n(x)\cos nx \\
        \psi_n(x)=\sin nx
    \end{cases}$

$\frac{a}{2}=\frac{1}{2}\frac{1}{\pi}\int_{-\pi}^{\pi}f(x)\,\mathrm{d}x=\frac{1}{\pi}\int_{-\pi}^{\pi}f(x)\frac{1}{\sqrt{2}}\,\mathrm{d}x\cdot\frac{1}{\sqrt{2}}$

\begin{remark}
    If $\left(V,\left<\cdot,\cdot\right>\right)$ is a finitely dimensional inner product space, then for ant orthonormal basis $\left\{v_1,\dots,v_n\right\}$, we have $v=\left<v,v_1\right>v_1+\dots+\left<v,v_n\right>v_n$
\end{remark}

\begin{theorem}[Dirichlet Convergence Theorem]
    Let $f$ be a piecewise $C^1$ function on $[-\pi,\pi]$, then the Fourier series of $f$ converges, and
    $$\frac{a_0}{2}+\sum_{n=1}^{\infty}\left(a_n\cos nx+b_n\sin nx\right)=\begin{cases}
            \frac{f(x+0)+f(x-0)}{2}      & x\in(-\pi,\pi) \\
            \frac{f(-\pi+0)+f(\pi-0)}{2} & x=\pm\pi
        \end{cases}$$
    Here, piecewise $C^1$ means $\exists t+0=-\pi<t_1<\dots<t_n=\pi$ s.t. $f$ is $C^1$ on $(t_i,t_{i+1})$, $i=0,1,2,\dots,n-1$ and $f_+'(t_i),f_-'(t_i)$ exists for $i=0,1,2,\dots,n$ and $f(t_i+0),f(t_i-0)$ exists
\end{theorem}

\begin{remark}
    For a function $f$ of period $2l$,
    $$g(x)=f\left(\frac{l}{\pi}x\right)$$

    $\frac{lT}{\pi}=2l\implies T=2\pi$

    Fourier series for $g$ is
    $$\frac{a_0}{2}+\sum_{n=1}^{\infty}\left(a_n\cos nx+b_n\sin nx\right),$$
    and the Fourier series for $f$ is
    $$\frac{a_0}{2}+\sum_{n=1}^{\infty}\left(a_n\cos\frac{n\pi}{l}x+b_n\sin\frac{n\pi}{l}x\right)$$
    where
    \begin{align*}
        a_n & =\frac{1}{\pi}\int_{-\pi}^{\pi}g(x)\cos nx\,\mathrm{d}x                                                           \\
            & =\frac{1}{\pi}\int_{-\pi}^{\pi}f\left(\frac{l}{\pi}x\right)\cos nx\,\mathrm{d}x                                   \\
            & \xlongequal{t=\frac{l}{\pi}x}\frac{1}{\pi}\int_{-l}^l f(t)\cdot\cos\frac{n\pi}{l}t\cdot\frac{\pi}{l}\,\mathrm{d}t \\
            & =\frac{1}{l}\int_{-l}^l f(x)\cos\frac{n\pi}{l}x\,\mathrm{d}x,                                                     \\
        b_n & =\frac{1}{l}\int_{-l}^l f(x)\sin\frac{n\pi}{l}x\,\mathrm{d}x
    \end{align*}
\end{remark}

\begin{example}
    $f(x)=x$

    The Fourier series is
    \begin{align*}
                                                         & \frac{\pi}{4}+\sum_{n=1}^{\infty}\left(\frac{(-1)^n-1}{n^2\pi}\cos nx+\frac{(-1)^{n+1}}{n}\sin nx\right) \\
        \xlongequal{\text{By Dirichlet Convergence Thm}} & \begin{cases}
                                                               f(x)          & x\in(-\pi,\pi) \\
                                                               \frac{\pi}{2} & x=\pm\pi
                                                           \end{cases}
    \end{align*}

    \begin{align*}
        a_0 & =\frac{1}{\pi}\int_{-\pi}^{\pi}f(x)\,\mathrm{d}x                                                                   \\
            & =\frac{1}{\pi}\int_0^{\pi}x\,\mathrm{d}x                                                                           \\
            & =\frac{\pi}{2}                                                                                                     \\
        a_n & =\frac{1}{\pi}\int_{-\pi}^{\pi}f(x)\cos nx\,\mathrm{d}x                                                            \\
            & =\frac{1}{\pi}\int_0^{\pi}x\cos nx\,\mathrm{d}x                                                                    \\
            & =\frac{1}{\pi}\left(\left.x\cdot\frac{\sin nx}{n}\right|_0^{\pi}-\int_0^{\pi}\frac{\sin nx}{n}\,\mathrm{d}x\right) \\
            & =\frac{1}{n^2\pi}\left.\cos nx\right|_0^{\pi}                                                                      \\
            & =\frac{(-1)^n-1}{n^2\pi}                                                                                           \\
        b_n & =\frac{1}{\pi}\int_{-\pi}^{\pi}f(x)\sin nx\,\mathrm{d}x                                                            \\
            & =\frac{1}{\pi}\int_0^{\pi}x\cdot\sin nx\,\mathrm{d}x                                                               \\
            & =\frac{1}{\pi}\left(\left.-\frac{x\cos nx}{n}\right|_0^{\pi}+\int_0^{\pi}\frac{\cos nx}{n}\,\mathrm{d}x\right)     \\
            & =-\frac{(-1)^n}{n}                                                                                                 \\
            & =\frac{(-1)^{n+1}}{n}
    \end{align*}
\end{example}

\begin{example}
    Let $a\notin\mathbb{Z}$, find the Fourier series of $f(x)=\cos ax$ on $[-\pi,\pi]$

    By translation, extends $f$ to a function $\tilde{f}$ on $\mathbb{R}$ with period $2\pi$

    \begin{align*}
        a_n
         & =\frac{1}{\pi}\int_{-\pi}^{\pi}\tilde{f}(x)\cos nx\,\mathrm{d}x                               \\
         & =\frac{2}{\pi}\int_0^{\pi}\cos ax\cdot\cos nx\,\mathrm{d}x                                    \\
         & =\frac{1}{\pi}\int_0^{\pi}\cos(a-n)x+\cos(a+n)x\,\mathrm{d}x                                  \\
         & =\frac{1}{\pi}\left.\left(\frac{\sin(a-n)x}{a-n}+\frac{\sin(a+n)x}{a+n}\right)\right|_0^{\pi} \\
         & =\frac{1}{\pi}\left(\frac{\sin(a-n)\pi}{a-n}+\frac{\sin(a+n)\pi}{a+n}\right)                  \\
         & =\frac{2a}{\pi}\frac{(-1)^n\sin a\pi}{a^2-n^2}                                                \\
        b_n
         & =0
    \end{align*}

    The Fourier series is
    $$\frac{a}{\pi}\frac{\sin a\pi}{a^2}+\sum_{n=1}^{\infty}\frac{2a}{\pi}\frac{(-1)^n\sin a\pi}{a^2-n^2}\cos nx\xlongequal{\text{Dirichlet Convergence Thm}}\cos ax,\quad x\in[-\pi,\pi]$$
\end{example}